\section{Contributions}
\label{introduction:contributions}

The benchmark places vector cloud-native formats and engines beside their traditional counterparts and runs both on identical Norwegian building data in the same cloud storage, addressing gaps that prior spatial benchmarks leave open through four contributions. The first, \acrshort{rq}1, compares the cloud-native combination of GeoParquet and DuckDB against PostGIS and Shapefile across spatial intersection, nearest-neighbor search, and bounding-box filtering. The second, \acrshort{rq}2, sets the distributed engine Apache Sedona against single-node PostGIS and DuckDB on a national-scale spatial join across small, medium, and large dataset tiers, a comparison no published benchmark has reported. This comparison identifies the scale at which a single node no longer suffices: at the large tier neither single-node engine completes the join while the distributed broadcast join does, so distribution there is a question of feasibility rather than speed. The third reports network transfer jointly with run-time and operational cost, a dimension absent from almost all spatial benchmarks. The fourth is methodological: it reports effect sizes, the magnitude of a difference rather than only its presence, and bootstrapped \acrlonglc{ci}s~(\acrshort{ci}). It also adopts a rank-based omnibus test, one overall test of whether the configurations differ, with post-hoc correction. These practices are rarely found in geospatial benchmarking.